\documentclass[uplatex,a4paper,11pt]{jsbook}
\usepackage[margin=22mm]{geometry}
\usepackage[dvipdfmx]{hyperref,xcolor}
\usepackage{pxjahyper}
\usepackage{fvextra,enumitem,longtable,booktabs,amssymb}
\usepackage{listings,jvlisting}

\hypersetup{
  colorlinks=true,
  linkcolor=blue!55!black,
  urlcolor=blue!65!black,
  pdftitle={LilyPond transpose Python版 利用・実装解説}
}
\DefineVerbatimEnvironment{command}{Verbatim}{
  fontsize=\small,frame=single,rulecolor=\color{gray!55},
  breaklines=true,breakanywhere=true
}
\setlist{nosep}
\lstset{
  language=Python,
  basicstyle=\ttfamily\scriptsize,
  numbers=left,
  numberstyle=\tiny,
  frame=single,
  breaklines=true,
  columns=fullflexible,
  keepspaces=true,
  showstringspaces=false,
  tabsize=4,
  captionpos=b
}
\lstdefinelanguage{Swift}{
  morekeywords={associatedtype,class,deinit,enum,extension,fileprivate,func,import,
    init,inout,internal,let,open,operator,private,protocol,public,rethrows,
    static,struct,subscript,typealias,var,break,case,continue,default,defer,
    do,else,fallthrough,for,guard,if,in,repeat,return,switch,where,while,as,
    catch,false,is,nil,super,self,Self,throw,throws,true,try},
  sensitive=true,
  morecomment=[l]{//},
  morecomment=[s]{/*}{*/},
  morestring=[b]"
}

\title{LilyPond\_transpose\\Python版・Swift版 利用／実装解説}
\author{smat1957}
\date{2026年9月8日}

\begin{document}
\maketitle

\begin{abstract}
本書は、LilyPondで記述された楽譜を指定音程だけ移調する
\texttt{LilyPond\_transpose}のPython版とSwift版を説明する。第I部ではPython版の
コマンドラインからの使い方と実装を、第II部ではSwift版のライブラリ、Apple各端末用
アプリ、および共通変換処理を、実際のプログラムソースを参照しながら解説する。
両実装について、字句解析、音楽理論計算、相対音高の解決、移調、LilyPond文字列への
復元という処理の対応関係も分かる構成とした。
本プログラムは完全なLilyPondパーサーではなく、バッハ無伴奏チェロ組曲をギター用に
移調・整形する実用範囲を主な対象としている。
\end{abstract}

\tableofcontents

\part{Python版}

\chapter{概要}

\section{目的}

Python版は、標準入力からLilyPondソースを読み、移調後のLilyPondソースを標準出力へ
書き出すコマンドラインプログラムである。入力ファイルを直接上書きしないため、変換前の
楽譜を残したまま別ファイルとして結果を生成できる。

主な処理対象は次のとおりである。

\begin{itemize}
  \item LilyPond音名と臨時記号
  \item \texttt{\textbackslash relative}ブロックとオクターブ記号
  \item 単音と和音\texttt{<...>}
  \item \texttt{\textbackslash key}による調指定
  \item 複数声部を含む並列ブロック
  \item コメント、音価、演奏記号など移調対象外の文字列の保持
\end{itemize}

\section{必要な環境}

移調処理にはPython 3を使用する。結果を譜面として確認するにはLilyPondも必要である。
リポジトリの代表的な構成は次のとおりである。

\begin{command}
LilyPond_transpose/
|-- python/       Python移調プログラム
|-- lilyp/        入力・出力LilyPondソース
|-- pdf/          生成したPDF
|-- swift/        Swift版
`-- gen*.sh       曲別の生成スクリプト
\end{command}

\chapter{基本操作}

\section{コマンド形式}

\begin{command}
cd python
python3 main.py SRC DST < input.ly > output.ly
\end{command}

\texttt{SRC}は移調元の基準音高、\texttt{DST}は移調先の基準音高である。単なる調名ではなく、
オクターブ指示も含む音高として解釈される。プログラムは

\begin{equation}
  \mathrm{shift}=\mathrm{pitch}(DST)-\mathrm{pitch}(SRC)
\end{equation}

によって半音移動量を求める。

\section{音名と臨時記号}

音名はLilyPond形式で指定する。代表例を表\ref{tab:pitch-name}に示す。

\begin{longtable}{@{}lll@{}}
\caption{代表的な音名}\label{tab:pitch-name}\\
\toprule
表記 & 意味 & 備考\\
\midrule
\endfirsthead
\toprule
表記 & 意味 & 備考\\
\midrule
\endhead
\texttt{c}, \texttt{d}, \texttt{e} & 自然音 & C, D, E\\
\texttt{cis}, \texttt{fis}, \texttt{gis} & シャープ & C\(\sharp\), F\(\sharp\), G\(\sharp\)\\
\texttt{des}, \texttt{ges} & フラット & D\(\flat\), G\(\flat\)\\
\texttt{es}, \texttt{as}, \texttt{bes} & フラット & E\(\flat\), A\(\flat\), B\(\flat\)\\
\bottomrule
\end{longtable}

通常用いる長調の主音として、\texttt{ces}, \texttt{ges}, \texttt{des}, \texttt{as},
\texttt{es}, \texttt{bes}, \texttt{f}, \texttt{c}, \texttt{g}, \texttt{d}, \texttt{a},
\texttt{e}, \texttt{b}, \texttt{fis}, \texttt{cis}を指定できる。入力中の
\texttt{\textbackslash major}または\texttt{\textbackslash minor}は維持されるため、
同じ引数形式を長調と短調の両方に使用する。

\section{オクターブ指示}

ASCIIアポストロフィー\texttt{'}は1個につき1オクターブ上、カンマ\texttt{,}は1個につき
1オクターブ下を表す。例えば次の二つは移調後の音名文字が同じでも方向が異なる。

\begin{longtable}{@{}lll@{}}
\toprule
指定 & 半音移動量 & 意味\\
\midrule
\texttt{g}から\texttt{d} & \(-5\) & 完全4度下\\
\texttt{g}から\texttt{d'} & \(+7\) & 完全5度上\\
\texttt{g}から\texttt{d,} & \(-17\) & 1オクターブと完全4度下\\
\texttt{c}から\texttt{c'} & \(+12\) & 1オクターブ上\\
\texttt{c}から\texttt{c,} & \(-12\) & 1オクターブ下\\
\bottomrule
\end{longtable}

bashやzshからアポストロフィーを渡す場合は、引数全体をダブルクォートで囲む方法が
分かりやすい。

\begin{command}
python3 main.py g "d'" < input.ly > output.ly
\end{command}

次のようにアポストロフィーだけをバックスラッシュでエスケープしても同じ結果になる。

\begin{command}
python3 main.py g d\' < input.ly > output.ly
\end{command}

バックスラッシュはシェルにアポストロフィーを解釈させないためのものであり、Pythonへ渡す
音高文字列の一部ではない。アプリやHTTP APIなど、シェルを介さず引数を渡す場合は
\texttt{d'}をそのまま使用する。アポストロフィーとカンマを混在させた\texttt{d',}は
受け付けない。

\chapter{使用例}

\section{C-DurからA-Durへ移調する}

\begin{command}
cd python
python3 main.py c a \
  < ../lilyp/cello/cello1009.ly \
  > ../lilyp/guitar/guitar1009A.ly
\end{command}

\texttt{c}から\texttt{a}は3半音下である。入力中の\texttt{\textbackslash key c
\textbackslash major}は\texttt{\textbackslash key a \textbackslash major}へ変換され、
音符と和音も同じ音程だけ移調される。

\section{C-DurからD-Durへ移調する}

\begin{command}
cd python
python3 main.py c d \
  < ../lilyp/cello/cello1009.ly \
  > ../lilyp/guitar/guitar1009D.ly
\end{command}

この指定は2半音上への移調である。

\section{G-DurからD-Durへ完全5度上げる}

\begin{command}
cd python
python3 main.py g "d'" \
  < ../lilyp/cello/cello1007.ly \
  > ../lilyp/guitar/guitar1007.ly
\end{command}

\texttt{d'}を指定することで、同じD-Durでも5半音下ではなく7半音上を選択する。

\section{変換結果をLilyPondで確認する}

移調後は必ずLilyPondでコンパイルし、構文、音高、調号、オクターブ、運指、譜面の読みやすさを
確認する。

\begin{command}
cd ../lilyp
lilypond BWV1009PreludeA.ly
mv BWV1009PreludeA.pdf ../pdf/
\end{command}

曲別の\texttt{gen1007.sh}、\texttt{gen1009A.sh}などを使えば、移調、PDF生成、既存結果との
比較をまとめて実行できる。

\chapter{対応範囲と注意事項}

本プログラムはLilyPondの全構文を扱う完全なパーサーではない。次の点に注意する。

\begin{itemize}
  \item 複雑なScheme式や独自マクロの内部は正しく解析できない場合がある。
  \item 移調元引数と入力中の調号が一致するかは自動確認しない。
  \item 移調後の綴りに通常範囲を超える臨時記号が必要な場合はエラーになることがある。
  \item \texttt{\textbackslash relative}と多声部の組合せは実際の楽譜で確認する必要がある。
  \item 変換結果を完成譜とせず、必ずLilyPondによるコンパイルと目視確認を行う。
\end{itemize}

\chapter{プログラムソースと処理内容}

\section{ソースコメントを手引きとして読む}

Python版のソースには、単なる処理内容だけでなく、LilyPond固有の相対音高規則、
和音後の基準音、異名同音の選択、構文をそのまま保存する理由などがコメントとして
丁寧に記録されている。このコメントは実装の補足ではなく、設計判断と注意点を伝える
重要な資料である。本章以降ではコード全体を掲載し、本文で処理の流れを俯瞰しながら、
細部についてはソースコメントと照合できるようにしている。

\chapter{全体構成}

\section{処理の流れ}

Python版は責務ごとに8ファイルへ分割されている。処理は次の順に進む。

\begin{command}
標準入力のLilyPond文字列
        |
        v
Tokenizer: 構文要素をToken列へ分解
        |
        v
Transpose: 音高と調号を移調
        |
        +-- RelativeResolver: 相対音高を絶対音高へ展開して再表記
        |
        v
Writer: Token列をLilyPond文字列へ復元
        |
        v
標準出力
\end{command}

\begin{longtable}{@{}p{36mm}p{105mm}@{}}
\toprule
ファイル & 主な責務\\
\midrule
\texttt{main.py} & 引数、標準入力、処理全体の制御\\
\texttt{tokens.py} & Tokenデータ型の定義\\
\texttt{tokenizer.py} & LilyPond文字列からToken列への変換\\
\texttt{music\_theory.py} & 音名、半音値、異名同音、綴りの計算\\
\texttt{pitch.py} & 音高文字列、オクターブ、MIDI値の相互変換\\
\texttt{relative.py} & LilyPond相対音高の解決と再生成\\
\texttt{transpose.py} & Token木の走査と移調処理の振り分け\\
\texttt{writer.py} & Token列からLilyPond文字列への復元\\
\bottomrule
\end{longtable}

\chapter{エントリーポイント}

\section{main.pyの役割}

\texttt{main.py}は二つの引数を検査し、\texttt{parse\_pitch()}で半音移動量を、
\texttt{letter\_shift\_between()}で音名文字の移動量を求める。この二種類の移動量を分ける
ことで、同じ半音値を持つ音でも、移調音程に沿った音名文字を選べる。

標準入力全体を\texttt{Tokenizer}へ渡し、\texttt{transpose\_tokens()}でToken列を変更した後、
各Tokenを\texttt{write\_token()}で標準出力へ書き出す。

\lstinputlisting[caption={エントリーポイント main.py}]{../python/main.py}

\chapter{Tokenによる中間表現}

\section{tokens.pyの役割}

文字列を直接置換すると、音名と変数名、コメント、コマンドなどを区別できない。そのため、
プログラムはいったん入力をToken列へ変換する。\texttt{NoteToken}は音名、オクターブ記号、
音価などの後続文字列を分離して保持する。\texttt{RelativeBlock}、\texttt{ChordToken}、
\texttt{ParallelBlock}は子Tokenを持ち、LilyPond構文の入れ子を表現する。

\lstinputlisting[caption={Token定義 tokens.py}]{../python/tokens.py}

\chapter{字句解析}

\section{tokenizer.pyの役割}

\texttt{Tokenizer}は入力位置を進めながら、コメント、コマンド、音符、和音、調指定、
\texttt{\textbackslash relative}、並列ブロックを読み取る。音名一覧は長い順に並べ、
\texttt{beses}を\texttt{b}として途中まで誤認しないようにしている。

音符は音名、\texttt{'}や\texttt{,}からなるオクターブ記号、空白までのsuffixへ分ける。
suffixを保持することで、音価、タイ、スラー、装飾などを移調後も書き戻せる。

\texttt{read\_relative()}は外側の波括弧に対応する範囲を切り出し、その内部を新しい
\texttt{Tokenizer}で再帰的に解析する。同様に和音や複数声部も子Token列として保持する。

\lstinputlisting[caption={字句解析 tokenizer.py}]{../python/tokenizer.py}

\chapter{音名と半音の理論計算}

\section{music\_theory.pyの役割}

\texttt{NOTE\_TO\_SEMITONE}は音名を0から11のピッチクラスへ変換する。
\texttt{transpose\_note\_name\_by\_interval()}は、半音移動量だけでなく音名文字の移動量も
受け取り、新しい文字と臨時記号を別々に決める。

例えばCからAへの移調では半音移動量は\(-3\)、文字移動量は5である。したがってCはA、
DはB、EはC\(\sharp\)というように、音程の文字関係を保った綴りが得られる。生成結果が
ダブルシャープまたはダブルフラットになる場合は、実用性を優先して代表的な綴りへ簡略化する。

\lstinputlisting[caption={音楽理論計算 music\_theory.py}]{../python/music_theory.py}

\chapter{絶対音高とオクターブ}

\section{pitch.pyの役割}

\texttt{split\_pitch()}は\texttt{fis,}を音名\texttt{fis}とオクターブ記号\texttt{,}へ分ける。
音名の後ろには\texttt{'}または\texttt{,}だけを許可し、上下記号の混在を拒否する。
\texttt{parse\_pitch()}はCをMIDI値60の基準とし、アポストロフィーの個数からカンマの個数を
引いた値に12を掛けてオクターブ差を加える。

\texttt{PitchPos}は音名、オクターブ、音名文字、MIDI値をまとめて保持する。
\texttt{midi\_to\_absolute\_lily()}は計算後のMIDI値を絶対LilyPond表記へ戻し、
\texttt{\textbackslash relative}の新しい基準音生成に使用される。

\lstinputlisting[caption={音高変換 pitch.py}]{../python/pitch.py}

\chapter{相対音高の解決}

\section{relative.pyの役割}

LilyPondの\texttt{\textbackslash relative}では、各音符の実際のオクターブが直前音との関係で
決まる。\texttt{RelativeResolver}は移調前と移調後の直前音をそれぞれ保持する。

単音を処理するときは、まず\texttt{resolve\_relative\_pitch()}で元の絶対音高を求め、半音移動量を
加え、移調後の直前音から必要な\texttt{'}または\texttt{,}を再計算する。これにより、音名だけを
置換する場合に起こる意図しないオクターブ移動を避ける。

和音では構成音を順に解決する一方、和音の次の単音へ引き継ぐ基準には和音の第1音を使う。
複数声部ではResolverを複製し、各声部が別の直前音状態を持つようにする。

\lstinputlisting[caption={相対音高処理 relative.py}]{../python/relative.py}

\chapter{Token列の移調}

\section{transpose.pyの役割}

\texttt{walk\_tokens()}はToken列を再帰的に走査し、種類に応じて処理を振り分ける。
\texttt{RelativeBlock}では専用Resolverを作り、通常の単音・和音では現在Resolverがあるか
どうかによって相対表記用または単純移調用の処理を選ぶ。

\texttt{KeyToken}はモードを維持したまま主音を移調する。並列ブロックではResolverを複製し、
一つの声部の終端音が別の声部の開始音へ影響しないようにする。

\lstinputlisting[caption={Token走査と移調 transpose.py}]{../python/transpose.py}

\chapter{LilyPond文字列への復元}

\section{writer.pyの役割}

\texttt{write\_token()}はTokenの種類ごとにLilyPond表記を組み立てる。
\texttt{RawToken}、\texttt{CommentToken}、\texttt{CommandToken}は内容を保持し、
\texttt{NoteToken}は音名、再計算済みオクターブ記号、suffixを連結する。入れ子を持つTokenは
子Tokenを再帰的に文字列化する。

TokenizerとWriterを対にすることで、移調対象以外の情報を可能な限り保存したまま、必要な
音高部分だけを書き換える。

\lstinputlisting[caption={LilyPond出力 writer.py}]{../python/writer.py}

\chapter{保守と検証}

\section{変更時に確認する事項}

音名や対応構文を追加するときは、一つのテーブルだけでなく、Tokenizer、音楽理論テーブル、
音高変換、Writerの整合性を確認する。少なくとも次の組合せを検証する。

\begin{itemize}
  \item 上方向、下方向、1オクターブをまたぐ移調
  \item シャープ系、フラット系、異名同音
  \item 単音、和音、調号
  \item \texttt{\textbackslash relative}内の上行・下行・明示的オクターブ指示
  \item 並列声部
  \item コメント、音価、スラー、タイなどの保持
  \item 不正な音名および\texttt{'}と\texttt{,}の混在の拒否
\end{itemize}

最終検証は文字列比較だけで終えず、生成したソースをLilyPondでコンパイルし、PDF上の音高と
譜面も確認する。

\part{Swift版}

\chapter{Swift版の目的と使い方}

\section{Python版との関係}

Swift版は、Python版と同じ移調処理をAppleプラットフォームから利用できるようにした実装で
ある。変換本体はSwift Packageの\texttt{LilyPondTransposeCore}へ集約し、macOS、iPad、
iPhoneの各SwiftUIアプリはこの共通ライブラリを呼び出す。これにより、画面構成やファイル
選択の方法が端末ごとに異なっても、音楽理論とLilyPond変換の処理は一か所で保守できる。

\section{動作環境とビルド}

\texttt{Package.swift}はSwift tools 6.0、macOS 14以降、iOS・iPadOS 17以降を指定する。
共通ライブラリとテストは次のコマンドでビルド、実行できる。

\begin{command}
cd swift
swift test
\end{command}

アプリは\texttt{LilyPondTranspose.xcodeproj}をXcodeで開き、macOS、iPad、iPhoneの
いずれかのSchemeと実行先を選ぶ。画面では移調元と移調先の音高を入力し、LilyPondソースを
読み込むか貼り付けて「変換」を実行する。macOSとiPadでは変換前後を横に並べ、画面幅の
限られるiPhoneではタブで切り替えて表示する。

\section{ライブラリからの利用}

公開APIは\texttt{LilyPondTransposer}の\texttt{transpose}メソッドである。

\begin{lstlisting}[language=Swift,caption={Swift Packageからの呼出し例}]
import LilyPondTransposeCore

let result = try LilyPondTransposer().transpose(
    lilyPondSource,
    from: "c",
    to: "a"
)
\end{lstlisting}

入力と出力はいずれも\texttt{String}であり、失敗時は\texttt{TransposeError}を送出する。
したがって、コマンドラインの標準入出力に依存せず、文書アプリや別のSwiftUI画面からも
同じ変換器を直接利用できる。

\chapter{Swift版の全体構成}

\begin{command}
SwiftUIアプリ (macOS / iPad / iPhone)
        |
        v
LilyPondTransposer: 公開API、入力の正規化
        |
        v
Tokenizer: LilyPond文字列をToken列へ分解
        |
        v
TransposeEngine: 音高、調号、相対音高を移調
        |
        v
LilyPondWriter: Token列を文字列へ復元
\end{command}

\begin{longtable}{@{}p{45mm}p{96mm}@{}}
\toprule
ファイル & 主な責務\\
\midrule
\texttt{LilyPondTransposer.swift} & 公開API、エラー型、入力音高の正規化\\
\texttt{Token.swift} & 列挙型による中間表現\\
\texttt{Tokenizer.swift} & LilyPond文字列の字句解析と入れ子構造の読取り\\
\texttt{MusicTheory.swift} & 音名、半音、異名同音、移調後の綴りの計算\\
\texttt{TransposeEngine.swift} & Token木の走査、相対音高の解決、移調\\
\texttt{LilyPondWriter.swift} & Token列からLilyPond文字列への復元\\
\texttt{*ContentView.swift} & 端末に適した入力、ファイル選択、結果表示\\
\texttt{LilyPondTransposer}\newline\texttt{Tests.swift} & 共通ライブラリの回帰テスト\\
\bottomrule
\end{longtable}

\lstinputlisting[language=Swift,caption={Swift Package定義 Package.swift}]{../swift/Package.swift}

\chapter{公開APIとエラー処理}

\texttt{LilyPondTransposer.transpose()}は、前後の基準音を正規化して半音移動量と音名文字の
移動量を計算し、Tokenizer、TransposeEngine、Writerを順番に接続する。UIで入力した
アポストロフィーがスマート引用符へ置換される場合があるため、\texttt{normalizePitch()}で
ASCIIの\texttt{'}へ戻してから検証する点は、Apple端末向け実装に固有の配慮である。

\texttt{TransposeError}は、不正な音高、閉じ括弧などを欠くソース、未対応の臨時記号を区別し、
\texttt{LocalizedError}として日本語の説明をUIへ渡す。

\lstinputlisting[language=Swift,caption={公開API LilyPondTransposer.swift}]{../swift/LilyPondTranspose/Core/LilyPondTransposer.swift}

\chapter{Tokenと字句解析}

\section{値型による中間表現}

Swift版の\texttt{Token}は\texttt{indirect enum}で定義する。関連値を使って、単音なら音名、
オクターブ、suffixを、和音や相対音高ブロックなら子Tokenを一つの型で表現する。
\texttt{Equatable}への準拠により、将来Token単位のテストも簡潔に記述できる。

\lstinputlisting[language=Swift,caption={Token定義 Token.swift}]{../swift/LilyPondTranspose/Core/Token.swift}

\section{Tokenizerの走査}

\texttt{Tokenizer}は文字列を\texttt{Character}配列として保持し、\texttt{position}を進める。
現在位置から\texttt{\textbackslash relative}、並列ブロック、コメント、調指定、コマンド、
和音、単音の順に判定する。音名は長い候補から照合するため、短い接頭辞による誤認を避ける。

波括弧の本体は深さを数えて切り出し、別のTokenizerで再帰的に解析する。閉じ記号がない場合は
その場で\texttt{malformedSource}を送出し、不完全なToken列を後段へ渡さない。

\lstinputlisting[language=Swift,caption={字句解析 Tokenizer.swift}]{../swift/LilyPondTranspose/Core/Tokenizer.swift}

\chapter{音楽理論と相対音高}

\section{MusicTheory}

\texttt{MusicTheory}は音名からピッチクラスへの表、自然音の半音位置、シャープ系・フラット系の
代表綴りを持つ。\texttt{transposeNote()}は半音差だけでなく音名文字の移動量も用いるため、
音程として自然な綴りを選べる。一方、絶対表記の簡易変換や調号には
\texttt{transposeSimple()}を使う。

Swiftの剰余演算は負数を返し得るため、\texttt{modulo()}で常に0以上へ正規化している。
これは下方向への移調で配列添字を安全に求めるために重要である。

\lstinputlisting[language=Swift,caption={音楽理論計算 MusicTheory.swift}]{../swift/LilyPondTranspose/Core/MusicTheory.swift}

\section{RelativeResolverとTransposeEngine}

\texttt{RelativeResolver}は移調前と移調後の直前音を\texttt{PitchPosition}として別々に保持する。
元の相対表記を絶対MIDI値へ解決し、移調後のMIDI値に対応する音名と、直前音から必要になる
アポストロフィーまたはカンマを再計算する。

和音では各構成音を順に処理するが、次の単音へ引き継ぐ基準は第1音である。並列ブロックでは
Resolverを声部ごとにコピーし、一方の声部の終端音が他方へ影響しないようにする。
\texttt{TransposeEngine.walk()}はこの状態を\texttt{inout}で受け渡しながらToken木を再帰走査する。

\lstinputlisting[language=Swift,caption={移調エンジン TransposeEngine.swift}]{../swift/LilyPondTranspose/Core/TransposeEngine.swift}

\chapter{LilyPond文字列への復元}

\texttt{LilyPondWriter}は列挙型の各caseを\texttt{switch}し、Tokenを文字列へ戻す。
コメント、コマンド、記号は保存した文字列をそのまま返し、変更対象の音符と調号だけを
組み立て直す。和音、相対音高、並列声部は子Tokenに対して\texttt{write()}を再帰適用する。

\lstinputlisting[language=Swift,caption={LilyPond出力 LilyPondWriter.swift}]{../swift/LilyPondTranspose/Core/LilyPondWriter.swift}

\chapter{Apple端末用ユーザーインターフェース}

三つのUIはいずれも\texttt{@State}に入力、結果、移調元、移調先、エラーを保持し、
\texttt{convert()}から共通の\texttt{LilyPondTransposer}を呼び出す。macOS版は
\texttt{NSOpenPanel}、iPad・iPhone版は\texttt{fileImporter}を利用する。iOS系では
セキュリティスコープ付きURLへのアクセスを開始し、\texttt{defer}で確実に終了する。

\lstinputlisting[language=Swift,caption={macOS画面 MacContentView.swift}]{../swift/LilyPondTranspose/macOS/MacContentView.swift}
\lstinputlisting[language=Swift,caption={iPad画面 PadContentView.swift}]{../swift/LilyPondTranspose/iPad/PadContentView.swift}
\lstinputlisting[language=Swift,caption={iPhone画面 PhoneContentView.swift}]{../swift/LilyPondTranspose/iPhone/PhoneContentView.swift}

\chapter{テストと保守}

Swift Testingによるテストは、単音、調号、和音、相対音高をまとめて移調する基本例に加え、
コメントとコマンドの保持、不正音名、移調方向を変えるオクターブ記号、スマート引用符の正規化、
上下記号の混在拒否を確認する。CoreをUIから分離したことで、シミュレータや画面操作なしに
変換ロジックを検証できる。

\lstinputlisting[language=Swift,caption={Swift版のテスト LilyPondTransposerTests.swift}]{../swift/LilyPondTransposeTests/LilyPondTransposerTests.swift}

Python版とSwift版の挙動を変更するときは、同じ入力に対する出力を相互比較し、最後にLilyPondで
コンパイルする。特に相対音高、和音直後の基準音、並列声部、異名同音は、文字列としてもっとも
らしくても意図した譜面と異なる可能性があるため、PDFの目視確認までを検証に含める。

\end{document}
